
This work presents EDGAR, a fine-grained Multimodal Retrieval-Augmented Generation (MRAG) architecture designed for querying educational documents in PDF format. The architecture combines document analysis, text extraction, and visual element description techniques to organize and retrieve information at the component level. This approach enables the use of textual and visual content to build the context for response generation while preserving the traceability of evidence to the original document. To demonstrate its functionality, an interactive web interface was implemented, allowing users to query documents and visualize the retrieved evidence. The architecture was subjected to a comparative experimental evaluation against traditional text-based RAG approaches, considering information retrieval and the quality of generated responses. The results highlight the contributions of visual enrichment to queries relying on figures and multimodal reasoning, as well as to the mean response quality score, although they do not indicate overall superiority over text-chunk indexing in initial page retrieval. In this sense, EDGAR aims to support more grounded and traceable interactions between students and educational materials, facilitating access to the textual and visual information contained in the documents.

\keywords{retrieval-augmented generation; multimodal retrieval; fine-grained retrieval; document processing; visual traceability.}
